% GENERATED FILE. Edit canonical lessons, then run npm run generate:books.
% canonical-source-hash: fnv1a64:f6d5dc31f54d3e4d
% canonical-lessons: RU-C34-joy, RU-C34-fear, RU-C34-anger, RU-C34-sadness, RU-C34-tiredness

\chapter{Joy, Fear, Anger, Sadness, Tiredness}
\label{ch:ru-joy-fear-anger-sadness-tiredness}
\hlchaptermodality{34}

\begin{chapteropening}
\textbf{By the end of this chapter:} \emph{I can name joy, fear, anger, sadness and tiredness.}

Complete the last lesson of chapter 34: I can name joy, fear, anger, sadness and tiredness.
\end{chapteropening}

\section[\texorpdfstring{\ru{радость} (rádost)}{радость (rádost)}]{\ru{радость} (rádost) — joy}

\label{lesson:RU-C34-joy}



\begin{quote}
Before the new one: say the Russian for hunger, then the Russian for thirst.
\end{quote}

\subsection*{You'll want to know: \ru{радость}}
\textbf{\ru{радость}} (\emph{rádost}) — \textquotedblleft{}joy\textquotedblright{}.

\begin{cousinweb}
\textbf{\ru{радость}} is \emph{\ru{рад}}, \textquotedblleft{}glad\textquotedblright{}, with \emph{-\ru{ость}}, the suffix that makes nouns of quality. \emph{\ru{рад}} goes back to Proto-Slavic \emph{*radŭ}; Vasmer finds a link with Old English \emph{rōt}, \textquotedblleft{}cheerful\textquotedblright{}, likely, but reports a rival account.
\end{cousinweb}

\begin{grammarlens}[title={What you've built}]
One more word for this chapter.
\end{grammarlens}

\subsection*{Your turn}
\begin{itemize}
\raggedright
  \item \emph{Say it:} \emph{rádost}
  \item \emph{Say it:} \emph{rádost}, once more
  \item \emph{Say it:} say \emph{rádost}
  \item \emph{Recall:} say the Russian for hunger, then the Russian for thirst, then say \emph{rádost} again
\end{itemize}

\begin{tcolorbox}[breakable,colback=teal!4,colframe=teal!35!black,title={Before you move on}]
What does \ru{радость} mean? (\textquotedblleft{}Joy\textquotedblright{}.) Say it once more.
\end{tcolorbox}
\section[\texorpdfstring{\ru{страх} (strakh)}{страх (strakh)}]{\ru{страх} (strakh) — fear}

\label{lesson:RU-C34-fear}



\begin{quote}
Before the new one: say the Russian for thirst, then the Russian for joy.
\end{quote}

\subsection*{You'll want to know: \ru{страх}}
\textbf{\ru{страх}} (\emph{strakh}) — \textquotedblleft{}fear\textquotedblright{}.

\begin{cousinweb}
\textbf{\ru{страх}} $\leftarrow$ Proto-Slavic \emph{*straxŭ} — inherited. Its deeper origin is not settled: one account Vasmer reports gives it a first sense of \textquotedblleft{}numbness\textquotedblright{} and compares Lithuanian \emph{stregti}, \textquotedblleft{}to freeze stiff\textquotedblright{}.
\end{cousinweb}

\begin{grammarlens}[title={What you've built}]
One more word for this chapter.
\end{grammarlens}

\subsection*{Your turn}
\begin{itemize}
\raggedright
  \item \emph{Say it:} \emph{strakh}
  \item \emph{Say it:} \emph{strakh}, once more
  \item \emph{Say it:} say \emph{strakh}
  \item \emph{Recall:} say the Russian for thirst, then the Russian for joy, then say \emph{strakh} again
\end{itemize}

\begin{tcolorbox}[breakable,colback=teal!4,colframe=teal!35!black,title={Before you move on}]
What does \ru{страх} mean? (\textquotedblleft{}Fear\textquotedblright{}.) Say it once more.
\end{tcolorbox}
\section[\texorpdfstring{\ru{гнев} (gnev)}{гнев (gnev)}]{\ru{гнев} (gnev) — anger}

\label{lesson:RU-C34-anger}



\begin{quote}
Before the new one: say the Russian for joy, then the Russian for fear.
\end{quote}

\subsection*{You'll want to know: \ru{гнев}}
\textbf{\ru{гнев}} (\emph{gnev}) — \textquotedblleft{}anger\textquotedblright{}.

\begin{cousinweb}
\textbf{\ru{гнев}} $\leftarrow$ Proto-Slavic \emph{*gněvŭ} — inherited. Vasmer thinks it probably belongs with \emph{\ru{гнить}}, \textquotedblleft{}to rot\textquotedblright{}, so the first sense would have been \textquotedblleft{}rot, poison\textquotedblright{}; others start from a sense \textquotedblleft{}to burn\textquotedblright{}.
\end{cousinweb}

\begin{grammarlens}[title={What you've built}]
One more word for this chapter.
\end{grammarlens}

\subsection*{Your turn}
\begin{itemize}
\raggedright
  \item \emph{Say it:} \emph{gnev}
  \item \emph{Say it:} \emph{gnev}, once more
  \item \emph{Say it:} say \emph{gnev}
  \item \emph{Recall:} say the Russian for joy, then the Russian for fear, then say \emph{gnev} again
\end{itemize}

\begin{tcolorbox}[breakable,colback=teal!4,colframe=teal!35!black,title={Before you move on}]
What does \ru{гнев} mean? (\textquotedblleft{}Anger\textquotedblright{}.) Say it once more.
\end{tcolorbox}
\section[\texorpdfstring{\ru{грусть} (grust)}{грусть (grust)}]{\ru{грусть} (grust) — sadness}

\label{lesson:RU-C34-sadness}



\begin{quote}
Before the new one: say the Russian for fear, then the Russian for anger.
\end{quote}

\subsection*{You'll want to know: \ru{грусть}}
\textbf{\ru{грусть}} (\emph{grust}) — \textquotedblleft{}sadness\textquotedblright{}.

\begin{cousinweb}
\textbf{\ru{грусть}} goes with the verb \emph{\ru{грустить}}, \textquotedblleft{}to be sad\textquotedblright{}. Vasmer reports the usual link with \emph{\ru{груда}}, \textquotedblleft{}a heap\textquotedblright{}, and with Lithuanian \emph{graudùs}, \textquotedblleft{}brittle; sad, moving\textquotedblright{}, from a root \textquotedblleft{}to crush\textquotedblright{} (Lithuanian \emph{grūsti}, \textquotedblleft{}to pound\textquotedblright{}).
\end{cousinweb}

\begin{grammarlens}[title={What you've built}]
One more word for this chapter.
\end{grammarlens}

\subsection*{Your turn}
\begin{itemize}
\raggedright
  \item \emph{Say it:} \emph{grust}
  \item \emph{Say it:} \emph{grust}, once more
  \item \emph{Say it:} say \emph{grust}
  \item \emph{Recall:} say the Russian for fear, then the Russian for anger, then say \emph{grust} again
\end{itemize}

\begin{tcolorbox}[breakable,colback=teal!4,colframe=teal!35!black,title={Before you move on}]
What does \ru{грусть} mean? (\textquotedblleft{}Sadness\textquotedblright{}.) Say it once more.
\end{tcolorbox}
\section[\texorpdfstring{\ru{усталость} (ustálost)}{усталость (ustálost)}]{\ru{усталость} (ustálost) — tiredness}

\label{lesson:RU-C34-tiredness}



\begin{quote}
Before the new one: say the Russian for anger, then the Russian for sadness.
\end{quote}

\subsection*{You'll want to know: \ru{усталость}}
\textbf{\ru{усталость}} (\emph{ustálost}) — \textquotedblleft{}tiredness\textquotedblright{}.

\begin{cousinweb}
\textbf{\ru{усталость}} is \emph{\ru{устать}}, \textquotedblleft{}to get tired\textquotedblright{}, with \emph{-\ru{ость}}, the suffix for nouns of quality. \emph{\ru{устать}} is \emph{\ru{у}-} with \emph{\ru{стать}}, \textquotedblleft{}to stand, to become\textquotedblright{} (Vasmer); Old Church Slavonic \emph{\ru{уставати}} meant \textquotedblleft{}to cease\textquotedblright{}: a tired person has come to a stop.
\end{cousinweb}

\begin{grammarlens}[title={What you've built}]
One more word for this chapter.
\end{grammarlens}

\subsection*{Your turn}
\begin{itemize}
\raggedright
  \item \emph{Say it:} \emph{ustálost}
  \item \emph{Say it:} \emph{ustálost}, once more
  \item \emph{Say it:} say \emph{ustálost}
  \item \emph{Recall:} say the Russian for anger, then the Russian for sadness, then say \emph{ustálost} again
\end{itemize}

\begin{tcolorbox}[breakable,colback=teal!4,colframe=teal!35!black,title={Before you move on}]
What does \ru{усталость} mean? (\textquotedblleft{}Tiredness\textquotedblright{}.) Say it once more.
\end{tcolorbox}
